\documentclass[11pt,a4paper]{article}
\usepackage[utf8]{inputenc}
\usepackage[T1]{fontenc}
\usepackage{amsmath,amssymb}
\usepackage{booktabs}
\usepackage{geometry}
\usepackage{enumitem}
\usepackage{xcolor}
\usepackage{hyperref}
\geometry{margin=2.4cm}
\definecolor{accent}{RGB}{38,70,83}
\hypersetup{colorlinks=true,linkcolor=accent,urlcolor=accent}

\title{\vspace{-1.5em}\textbf{Document 2 --- The WILD Dataset}\\
\large Collection, Representation, Sessions, and Splits}
\author{}
\date{}

\begin{document}
\maketitle
\tableofcontents
\vspace{1em}\hrule\vspace{1em}

%% ==================================================================
\section{Overview}

All experiments use the \textbf{WILD} dataset (\emph{WiFi-based Indoor Localization Dataset}),
the public data released with DLoc by UC~San~Diego. WILD was collected with \textbf{MapFind}, an
autonomous data-collection robot, and provides, for each spatial point, the CSI heard from every
AP \emph{together with} a high-accuracy ground-truth position. This combination of fine-grained
wireless measurements and reliable labels is what makes supervised deep localization possible.

%% ==================================================================
\section{How the data was collected: MapFind}

MapFind is a TurtleBot2 carrying: a Hokuyo \textbf{LIDAR}, an Orbbec \textbf{RGB-D camera}, and a
\textbf{Quantenna WiFi card} at phone height. As it drives the space at a constant
$15$\,cm/s (slow enough to avoid Doppler effects), it simultaneously:
\begin{itemize}[nosep]
  \item maps the environment and localizes itself with SLAM (\textbf{RTAB-Map}, median ground-truth
  error $\approx 5.7$\,cm --- validated against an HTC~Vive VR tracker), and
  \item records timestamped CSI from every AP (each AP estimates the channel every $50$\,ms).
\end{itemize}
The CSI plus the SLAM position form one labelled training example. About $12{,}500$ points are
collected per scenario; the full release spans $\sim$100k points.

%% ==================================================================
\section{Two environments, eight scenarios}

WILD covers two physical spaces:

\begin{table}[h]
\centering
\caption{The two WILD environments. This project uses the \textbf{complex (Jacobs Hall)} space.}
\begin{tabular}{@{}lccl@{}}
\toprule
Environment & Area & APs & Scenarios \\
\midrule
Simple (Atkinson Hall)  & 500\,sq.\ ft.  & 3 & 3 (two days + one reflector) \\
\textbf{Complex (Jacobs Hall)} & 1500\,sq.\ ft. & \textbf{4} & \textbf{5} (two times + furniture/reflector) \\
\bottomrule
\end{tabular}
\end{table}

The complex space is harder: AP~4 is hidden behind a wall (strong NLOS) and a wall of plasma
screens behind AP~3 creates rich multipath. \textbf{All results in this project use the four-AP
Jacobs Hall data}; the Atkinson (3-AP) space is not used.

%% ==================================================================
\section{The five Jacobs Hall sessions}

\begin{table}[h]
\centering
\caption{The five Jacobs Hall collection sessions used throughout the project.}
\label{tab:sessions}
\begin{tabular}{@{}llll@{}}
\toprule
\# & Session (file stem) & Condition & Samples \\
\midrule
1 & \texttt{jacobs\_July28}        & standard, no furniture            & (see below) \\
2 & \texttt{jacobs\_July28\_2}     & same setup, $\sim$1\,h later      & 17{,}574 train$^\dagger$ / 4{,}260 test \\
3 & \texttt{jacobs\_Aug16\_1}      & random furniture added            & 11{,}201 \\
4 & \texttt{jacobs\_Aug16\_3}      & different furniture arrangement   & 8{,}754 \\
5 & \texttt{jacobs\_Aug16\_4\_ref} & furniture $+$ aluminium reflector & 5{,}486 \\
\bottomrule
\end{tabular}
\end{table}
{\footnotesize $^\dagger$ The standard ``Fig.\,10b'' split combines Session~1 and the training
part of Session~2 into 17{,}574 training samples; the held-out Session~2 test set has 4{,}260
samples.}

Sessions 3--5 are captured three weeks later under progressively harder conditions and are used
to measure \textbf{cross-session generalization} (does a model trained earlier still work after
the room changes?). The reflector in Session~5 introduces strong specular multipath and is the
hardest case.

%% ==================================================================
\section{Feature representation}

Each training example is a stack of per-AP 2D heatmaps over the floorplan:
\[
\text{input} \in \mathbb{R}^{\,4 \times 161 \times 361}, \qquad
\text{label} \in \mathbb{R}^{\,1 \times 161 \times 361},
\]
with a spatial resolution of \textbf{0.1\,m/pixel} spanning $16.1\,\text{m}\times36.1\,\text{m}$.
Three arrays are stored per \texttt{.mat} file:

\begin{table}[h]
\centering
\caption{Variables stored in each feature \texttt{.mat} file.}
\begin{tabular}{@{}lll@{}}
\toprule
Variable & Shape & Role \\
\midrule
\texttt{features\_w\_offset}  & $[N,4,161,361]$ & network \textbf{input} (heatmaps with ToF offset) \\
\texttt{features\_wo\_offset} & $[N,4,161,361]$ & \textbf{target} for the consistency decoder (offset removed) \\
\texttt{labels\_gaussian\_2d} & $[N,161,361]$   & \textbf{target} for the location decoder (Gaussian at GT) \\
\bottomrule
\end{tabular}
\end{table}

\paragraph{From CSI to heatmaps.} The MATLAB pipeline (\texttt{CSI\_to\_DLocFeatures}) converts
each raw complex channel matrix into a 2D AoA--ToF multipath profile and projects it onto the
physical grid:
\begin{itemize}[nosep]
  \item distance search grid \texttt{D\_VALS} $= -10\!:\!0.1\!:\!30$\,m,
  \item angle search grid \texttt{THETA\_VALS} $= -\pi/2\!:\!0.01\!:\!\pi/2$,
  \item the ground-truth label is rendered as a 2D Gaussian of spread $\sigma = 0.3$\,m at the
  SLAM position.
\end{itemize}
The ``with-offset'' and ``without-offset'' versions differ only in whether the per-AP ToF offset
has been compensated.

%% ==================================================================
\section{File format note (for reproducers)}

The \texttt{.mat} files are HDF5 (MATLAB v7.3). HDF5 stores MATLAB dimensions \emph{reversed}, so
a MATLAB array $[N,4,161,361]$ appears in Python/h5py as $[361,161,4,N]$ and must be transposed
on load (the data loader handles this). Files are large (the Session~3 features are $\sim$15\,GB)
because they store dense per-AP heatmaps for every sample.

%% ==================================================================
\section{How the splits are used in this project}

\begin{description}[nosep,leftmargin=2.0cm,style=nextline]
  \item[Standard split (in-distribution)] Train on Session~1 $+$ Session~2-train
  ($17{,}574$ samples); test on Session~2-test ($4{,}260$). Used for accuracy, robustness,
  scheduler, and multi-seed experiments.
  \item[Initial cross-session (2-session)] Train only on the two July sessions; test on
  Sessions 3--5. Simpler but \emph{not} comparable to the DLoc paper.
  \item[Official cross-session (leave-one-session-out)] Train on three setups (the two July
  sessions $+$ two of the three August sessions) and test on the held-out August session ---
  identical to the DLoc paper's Table~1. Used for the final, fair comparison.
\end{description}

\paragraph{Hosting.} The processed feature \texttt{.mat} files and the trained weights are
archived on HuggingFace (\texttt{JustAGeek/dloc-wild-fig10b} for data,
\texttt{JustAGeek/dloc-code} for code/weights/results); raw channel matrices are kept locally.
The experiments that consume these splits are described in \texttt{04\_experiments}.

\end{document}
